\section*{Chapitre \ref{ch:g3:food-chains} --- Les chaînes alimentaires}

\begin{solution}{exo:g3:food-chains:1}
«~\omterm{def:g3:food-chains:chain}{Est mangé par}~». Elle va du mangé vers le mangeur --- elle suit la
nourriture.
\end{solution}

\begin{solution}{exo:g3:food-chains:2}
Une \omterm{def:g1:plants-around-us:plant}{plante} (ou un autre \omterm{def:g1:living-or-not:living}{être vivant} qui fabrique sa propre
nourriture). Parce que seules les \omterm{def:g1:plants-around-us:plant}{plantes} se nourrissent sans manger
personne --- à partir de la lumière, de l'eau et de l'air~; chaque
maillon \omterm{def:g1:animals-around-us:animal}{animal} vit de ce premier maillon vert, directement ou par
d'autres.
\end{solution}

\begin{solution}{exo:g3:food-chains:3}
Herbe $\rightarrow$ sauterelle $\rightarrow$ lézard $\rightarrow$
buse.
\end{solution}

\begin{solution}{exo:g3:food-chains:4}
Les flèches sont à l'envers~: herbe $\rightarrow$ lapin
$\rightarrow$ renard --- l'herbe est mangée par le lapin, qui \omterm{def:g3:food-chains:chain}{est
mangé par} le renard.
\end{solution}

\begin{solution}{exo:g3:food-chains:5}
\omterm{def:g1:plants-around-us:parts}{Feuilles} de chêne $\rightarrow$ \omterm{ex:g2:animals-grow-up:butterfly}{chenille} $\rightarrow$ mésange bleue
$\rightarrow$ épervier.
\end{solution}

\begin{solution}{exo:g3:food-chains:6}
La salade est mangée par l'escargot~; le hérisson mange l'escargot.
\end{solution}

\begin{solution}{exo:g3:food-chains:7}
Par exemple~: \omterm{def:g1:plants-around-us:plant}{plantes} aquatiques $\rightarrow$ têtard
$\rightarrow$ épinoche $\rightarrow$ héron.
\end{solution}

\begin{solution}{exo:g3:food-chains:8}
Parce que la nourriture de la nourriture de sa nourriture est
l'herbe~: moins de sauterelles sur l'herbe grillée veut dire moins
de lézards, et moins de lézards veut dire une chasse maigre pour la
buse. Un changement dans le premier maillon remonte toute la chaîne.
\end{solution}

\begin{solution}{exo:g3:food-chains:9}
Empoisonner les escargots affame leur mangeur~: le hérisson s'en va
ou meurt. Plus tard, le hérisson --- principal ennemi des escargots
--- ayant disparu, les limaces et les escargots survivants se
multiplient plus fort qu'avant~: enlever un maillon voyage à la fois
vers le haut de la chaîne (hérisson affamé) et vers le bas
(escargots non chassés).
\end{solution}

\begin{solution}{exo:g3:food-chains:10}
Vrai --- par le milieu de la chaîne. Le lézard mange des
sauterelles, et les sauterelles mangent de l'herbe~: si l'herbe
manque, le manque remonte, maillon après maillon, de l'herbe à la
sauterelle puis au lézard. Le lézard ne touche jamais à l'herbe, et
il en dépend tout de même.
\end{solution}
